% 第六节：无穷远
\chapter{无穷远}

\section{无穷远点与无穷远直线}

由齐次坐标一节可知，分量形如 $(X,Y,0)$ 的点矢不能对应到平面上的点，因为这会使得 $X/Z$、$Y/Z$ 的分母为零。为使这样的点矢也有对应的几何对象，给平面添加一些新的元素，称为无穷远点。

由点矢与线矢一节又可知，形如 $(0,0,C)$ 的线矢不对应任何直线，因为它给出的直线方程形如 $0\cdot x+0\cdot y+C=0$。为使这样的线矢也有对应的几何对象，再给平面添加一个新的元素。

\begin{definition}[无穷远直线]
线矢 $(0,0,C)$（$C\neq 0$）所对应的直线称为无穷远直线，记作 $\linf$。在不引起混淆时也简写作 $\infty$。
\end{definition}

\begin{definition}[射影平面]
平面上的全体点，连同上面添加的全体无穷远点，合起来称为（实）射影平面。全体无穷远点构成无穷远直线 $\linf$。
\end{definition}

\section{平行线的交点}

在点矢与线矢一节中已经知道，两条不平行的直线有唯一的交点。引入无穷远点之后，平行的直线也有交点：它们的交点是无穷远点。下面的命题给出这一交点的齐次坐标。

\begin{proposition}\label{prop:parallel}
设两条直线的线矢为
\[
m_i=(A,B,C_1),\qquad n_i=(A,B,C_2),\qquad C_1\neq C_2,
\]
则两条直线平行，它们的交点为
\[
m\times n=\big(B(C_2-C_1),\,A(C_1-C_2),\,0\big)=(C_2-C_1)\,(B,-A,0),
\]
这是一个无穷远点。
\end{proposition}

\begin{proof}
两条直线的方向由 $(A,B)$ 决定，$m,n$ 的方向相同，故它们平行。直接按 $(A\times B)_i=\epsilon_{ijk}A^jB^k$ 计算：
\[
(m\times n)_1=m^2n^3-m^3n^2=B C_2-C_1B=B(C_2-C_1),
\]
\[
(m\times n)_2=m^3n^1-m^1n^3=C_1A-AC_2=A(C_1-C_2),
\]
\[
(m\times n)_3=m^1n^2-m^2n^1=AB-BA=0 .
\]
其第三个分量为零，故为无穷远点。
\end{proof}

\begin{remark}
两个分量的符号搭配容易记错，核对的方法是：直线 $Ax+By+C=0$ 的方向为 $(-B,A)$，因此上面
的无穷远点应当与 $(B,-A,0)$ 成比例，计算结果表明确实如此。同时又可见两条平行线
（$C_1\neq C_2$）交于\emph{同一个}无穷远点，因为该点只依赖于 $(A,B)$ 而不依赖于 $C$。
\end{remark}

下面的结论说明无穷远点与无穷远直线在射影平面中的基本性质。

\begin{proposition}
在射影平面中有下列结论：
\begin{enumerate}
\item 两条平行线交于同一个无穷远点；
\item 无穷远直线与任意一条直线都交于一个无穷远点；
\item 无穷远点在无穷远直线上；
\item 每条直线都经过一个无穷远点，不平行的直线经过不同的无穷远点；
\item 过两个无穷远点所作的直线是无穷远直线；
\item 设 $m$ 是一条直线，则当 $\lambda$ 取遍全体实数时，$m+\lambda\linf$ 给出所有与 $m$ 平行的直线。
\end{enumerate}
\end{proposition}

\begin{proof}
逐条验证。设直线 $m_i=(A,B,C_1)$、$n_i=(A',B',C_2)$。

(1) 平行意味着 $(A',B')=k(A,B)$，此时
\begin{align*}
(m\times n)_1&=B C_2-C_1B'=B(C_2-kC_1),\\
(m\times n)_2&=C_1A'-AC_2=A(kC_1-C_2),\\
(m\times n)_3&=AB'-A'B=0,
\end{align*}
其方向与 $(B,-A,0)$ 相同。凡方向与 $(A,B)$ 相同的直线都得到同一个点 $(B,-A,0)$，
故它们交于同一个无穷远点。

(2) 取 $\linf$ 的代表元 $(0,0,1)$，则
\[
\linf\times m=(0,0,1)\times(A,B,C)=(-B,\,A,\,0),
\]
第三个分量为零，是无穷远点。

(3) 无穷远点形如 $(X,Y,0)$，与 $\linf=(0,0,1)$ 缩并得 $0\cdot X+0\cdot Y+1\cdot 0=0$，故它在 $\linf$ 上。

(4) 由 (2)，直线 $m$ 与 $\linf$ 的交点为 $(-B,A,0)$，它只由方向 $(A,B)$ 决定；不平行的两条直线方向不成比例，所得的点也不同。

(5) 设两个不同的无穷远点为 $p=(X_1,Y_1,0)$、$q=(X_2,Y_2,0)$，则
\[
p\times q=(0,\,0,\,X_1Y_2-Y_1X_2),
\]
其第三个分量非零，故所得线矢是 $\linf$，即过 $p,q$ 的直线是无穷远直线。

(6) 取 $m_i=(A,B,C)$、$\linf=(0,0,1)$，则 $m+\lambda\linf=(A,B,C+\lambda)$。这些直线的方向都是 $(A,B)$，故都与 $m$ 平行；反过来，任何与 $m$ 平行的直线都可写成 $(A,B,C')=m+(C'-C)\linf$。故 $m+\lambda\linf$ 恰好给出所有与 $m$ 平行的直线。
\end{proof}

\section{无穷远结构与几何的层次}

无穷远直线的定义依赖于坐标：它的前两个分量被取为零。这一条件不可能写成与坐标系无关的形式，因为换一个坐标系之后，它的前两个分量一般不再为零。因此，齐次坐标一节中“抹去坐标系的一切痕迹”这第三步对无穷远结构不能实现，无穷远结构并不完全属于射影几何。

事实上，无穷远结构是三维矢量空间上的一个附加结构，只需要如下定义就可以产生。

\begin{definition}[仿射平面]
在三维矢量空间的全体线矢中指定一个特殊的线矢，称为无穷远线矢，记作 $\linf$，就在相应的射影几何中产生了一个无穷远结构。带无穷远结构的射影平面称为仿射平面。在此基础上可以进一步定义：无穷远点就是位于无穷远直线上的点；两条直线平行就是它们相交于无穷远点。
\end{definition}

于是，几何的层次关系可以更清楚地看到：

\begin{itemize}
\item 射影几何处于最底层，它只涉及点矢、线矢以及坐标系无关的运算，例如判断若干点是否共线、若干直线是否共点；
\item 仿射几何在射影几何之上多出了无穷远结构，因而涉及更多的关系，例如两条直线是否平行；
\item 解析几何处于最上层，它在仿射几何之上又多出了坐标系结构，因而涉及最多的关系，例如平行、垂直、长度，乃至 $y=\sin x$ 这样的曲线。
\end{itemize}

层次越高，可以表达的关系越多；高层次的关系可以表示低层次的关系，但低层次的关系未必能表示高层次的关系。
